\section{The \phenaki model}
\label{sec:method}
Inspired by the previous work in auto-regressive text to image~\cite{DALLE,PARTI,saharia2022photorealistic} and text to video~\cite{nuwa,wu2022nuwa,cogvideo}, \model is designed with two main components (see Figure~\ref{fig:models}): an encoder-decoder model which compresses videos to discrete embeddings (i.e. tokens) and a transformer model to \textit{translate} text embeddings to video tokens. To get the text embeddings, \model uses a pre-trained language model, T5X~\cite{roberts2022scaling}. We will discuss each one of these components in the following subsections.
 
% \subsection{Text embeddings}
% \rubville{Depending on results, talk about either learned text encoder or pre-computed t5x embeddings}
% \begin{itemize}
%     \item Reference to this model?
%     \item Sequence length is 64 tokens
%     \item embedding size is 4096
%     \item motivate using a pre-trained text tokenizer that it saves us memory since we can pre-compute embeddings and we do not need to load it into memory. This frees up space for compute. Imagen also had good results with it.
% \end{itemize}
%and trained key components that that work together to generate video from a text prompt. Additionally there is also a pre-trained text tokenizer based on T5X.
%The first component is the video (de-)tokenizer. It is able to take a video and return a compressed tokenized representation of this video. 
%text to video. The first model is an encoder-decoder transformer that quantizes video into sp89atio-temporal tokens. The second model maps textual descriptions into the spatio-temporal tokens computed using the encoder/decoder model. The tokens generated by the second model are then passed to the decoder in the first model to generate video. Next, we describe the details of each of our models, training strategies to learn from image and video datasets, super-resolution models, and other contributions that enable high-performing open domain text-to-video generation.

\subsection{Encoder-decoder video model: \cvivit}
% \rubville{Mention difference between the original vivit architecture and our non-causal \vivit architecture. Make sure reviewers don't think it's exactly the same.}
One of the primary challenges for generating video from text, is to get a compressed representation of videos. Previous work on text to video either use per-frame image encoders~\cite{cogvideo,nuwa,harp} such as VQ-GAN~\cite{vqgan} or fixed length video encoders~\cite{godiva} such as VideoVQVAE~\cite{videovqvae}. The former allows for generating videos of arbitrary length, however in practice, the videos have to be short because the encoder does not compress the videos in time and the tokens are highly redundant in consecutive frames. The latter is more efficient in the number of tokens but it does not allow to generate variable length videos. In \model, our goal is to generate videos of variable length while keeping the number of video tokens to a minimum so they can be modeled with a transformer within current computational limitations. To do so, we introduce \cvivit, a causal variation of \vivit~\cite{vivit} with additional architectural changes for video generation, which can compress the videos in temporal and spatial dimensions, while staying auto-regressive in time, This capability allows for generating videos of arbitrary length auto-regressively.
% given the reduuse a per frame image encoder to get video, who relied on a per frame model, as is typically used in text to image research. \pikinder{Can someone get us an illustration of this this}\rubville{It is difficult to notice the differences in a pdf. The FVD metric should be enough to show this}. However this approach has 2 key issues. First it does not use the redundancy between frames to encode information effectively. Second, when frames are decoded independently it can introduce temporal artefacts. We will discuss our basic ViVit encoder-decoder model first and discuss its limitations for autoregressive extension of video, which prevents it from being used in story mode. Then we will introduce C-ViVit that re-enables autogressive video generation.

\paragraph{Encoder architecture:}
As illustrated in Figure \ref{fig:models}, we start with a video sequence of $t_x + 1$ frames with a resolution of $w_x\times h_x$ and $c_x$ channels: $\bx \in \mathds{R}^{(t_x + 1) \times h_x \times w_x \times c_x}$. This sequence will be compressed into a token representation of size $(t_z + 1)\times w_z \times h_z$ where the first $w_z\times h_z$ tokens represent the first frame independently from the rest of the video, and the remaining tokens represent spatio-temporal video tokens that auto-regressively depend on previous frames. 
To do so, we extract non-overlapping image patches of size $w_p \times h_p \times c_p$ from the first frame and video patches of size $t_p\times w_p \times h_p \times c_p$ from the rest of the video.
We typically use all channels at once such that the number of patches equals the number of video tokens $t_z=\frac{t_x}{t_p}$, $w_z=\frac{w_x}{w_p}$ and  $h_z=\frac{h_x}{h_p}$.
Each of these patches is flattened and linearly projected into a $d_z$ dimensional space.
We combine the spatial dimensions to have a tensor of shape $(t_z + 1)\times w_z*h_z \times d_z$ where the spatial and temporal dimensions are separated.
Then multiple transformer layers are applied along the spatial dimensions with all-to-all attention.
This is followed by multiple transformer layers over the temporal dimension with causal attention such that each spatial token only observes spatial tokens from previous frames in an auto-regressive manner.
The effect of this is that the first frame can be completely independently encoded.
This opens up the possibility of text to image training to be embedded naturally into our video model.
The second advantage is that we can condition the video generation process on a number of starting frames. 
The resulting patch embeddings $\bz$ of shape $t_z\times w_z \times h_z \times d_z$ are then tokenized into learned codewords $c_z$ by vector quantization.
The codebook learning will be discussed later together with the losses.
% \pikinder{Expand this to explain the genius and the novelty of this}

\paragraph{Decoder architecture:}
The \cvivit decoder is simply an upside down version of the encoder. First tokens are transformed into embeddings. This is followed by the temporal transformer, then the spatial transformer.
After the output of the spatial transformer, we apply a single linear projection without activation to map the tokens back to pixel space.

\paragraph{Quantization and Losses:}
To learn a discrete latent space, we quantize our encoder outputs into the entries of a learned codebook via the vector quantization (VQ) objective in VQVAEs \citep{vqvae},
\begin{equation}
    L_{\textit{VQ}} = \| \text{sg}(\bz) - \be \|^2_2 + \beta \| \bz - sg(\be) \|^2_2,
\end{equation}
where $\text{sg}(x) \equiv x$, and $\frac{\text{d}}{\text{d}x} \text{sg}(x) \equiv 0$ is the stop-gradient operator, $\beta$ is the commitment loss weight, and $\be$ is a codebook vector from codebook $\mathbf{E}$. The index to the codebook vector closest to $\bz$ is found by $i = \text{argmin}_j \| \bz - \mathbf{E}_j \|^2_2$. 
In addition to the VQ objective, we adopt the factorized and $\ell_2$-normalized codes from ViT-VQGAN~\citep{vit_vqgan} to improve codebook usage and reconstruction quality.

To train our model, we use a combination of $L_2$ loss, image perceptual loss $L_{\textit{IP}}$~\citep{JohnsonPerceptual2016,ZhangPerceptual2018}, video perceptual loss $L_{\textit{VP}}$ by using the I3D network~\citep{i3d} as feature extractor, and adversarial loss $L_{\textit{Adv}}$ with StyleGAN architecture~\citep{stylegan}.
As training objective, we use the following
\begin{equation}
    L = L_{\textit{VQ}} + 0.1\times L_{\textit{Adv}} + 0.1\times L_{\textit{IP}} + 1.0\times L_{\textit{VP}} + 1.0\times L_2.
\end{equation}

\paragraph{Novelty over the \vivit architecture:}
While our proposed \cvivit architecture is inspired by the factorized encoder in \vivit~\cite{vivit}, we modify their architecture to enable self-supervised learning from unlabeled videos. We first remove the \texttt{[CLS]} tokens in the spatial and the temporal transformers. Next, we apply temporal transformer for all spatial tokens computed by the spatial encoder, in contrast to  single run of the temporal transformer over the \texttt{[CLS]} tokens in \vivit.
% We denote the modification of their factorized encoder as \vivit encoder.
Most importantly, the \vivit encoder requires a fixed length video input due to the all-to-all attention in time. Therefore, we apply causal attention instead such that our \cvivit encoder becomes auto-regressive and allows for a variable number of input frames which are necessary to learn from image datasets, and auto-regressively extrapolate video or single frames into the future. 


% \subsubsection{The \vivit encoder-decoder architecture}
% Our encoder model allows to generate a compressed representation of the video in the form of video tokens. Our decoder can reconstruct a video based on the tokens. The ViViT encoder-decoder model processes video as a single unit, not per frame as in previous works. 

% \paragraph{ViViT Encoder:}
% We start with a video sequence of $V_T$ frames with a resolution of $V_W\times V_H$ and $V_C$ channels: $\bx \in \mathds{R}^{V_T \times V_H \times V_W \times V_C}$. This sequence will be compressed into a token representation of size $Z_T\times Z_W\times Z_H$. 
% To do so, first we extract non-overlapping patches of size $P_T\times P_W \times P_H \times P_C$. We typically use all channels at once such that the number of patches equals the number of video tokens $Z_T=\frac{V_T}{P_T}$, $Z_W=\frac{V_W}{P_W}$ and  $Z_H=\frac{V_H}{P_H}$. Each of these patches is flattened and linearly projected into a $Z_D$ dimensional space.
% We combine the spatial dimensions to have a tensor of shape $Z_T\times Z_W*Z_H \times Z_D$ where the spatial and temporal dimensions are separated.
% Then multiple transformer layers are applied along the spatial dimensions (each frame is processed independently). This is followed by multiple transformer layers over the temporal dimension (each location is processed independently). The resulting patch embeddings $\bz$ of shape $Z_T\times Z_W \times Z_H \times Z_D$ are then tokenized into learned codewords $Z_C$ by vector quantization. The codebook learning will be discussed later together with the training losses.
% % Each of these these patch embeddings need to be discretized to obtain the tokenized representation $\bz\in$ need to be discretized into the $Z_C$ codewords. This will be discussed later together with the losses.

% \paragraph{ViViT Decoder:}
% The ViViT decoder is simply an upside down version of the encoder. First tokens are transformed into embeddings. This is followed by the temporal transformer, then the spatial transformer.
% After the output of the spatial transformer, we apply a single linear projection without activation to map the tokens back to pixel space.
% % After the output of the spatial transformer, we apply a single 3D transpose convolution without activation to map the tokens back to pixel space.

% \paragraph{ViViT Quantizer and losses:}
% To learn a discrete latent space, we quantize our encoder outputs into the entries of a learned codebook via the vector quantization (VQ) objective in VQVAEs \citep{vqvae},
% \begin{equation}
%     L_{VQ} = \| \text{sg}(\bz) - \be \|^2_2 + \beta \| \bz - sg(\be) \|^2_2,
% \end{equation}
% where $\text{sg}(x) \equiv x$, and $\frac{\text{d}}{\text{d}x} \text{sg}(x) \equiv 0$ is the stop-gradient operator, $\beta$ is the commitment loss weight, and $\be$ is a codebook vector from codebook $\mathbf{E}$. The index to the codebook vector closest to $\bz$ is found by $i = \text{argmin}_j \| \bz - \mathbf{E}_j \|^2_2$. 
% In addition to the VQ objective, we adopt the factorized and $\ell_2$-normalized codes from ViT-VQGAN \citep{vit_vqgan} to improve codebook usage and reconstruction quality.

% To train our model, we use a combination of $l2$ loss, image perceptual loss \citep{JohnsonPerceptual2016,ZhangPerceptual2018}, video perceptual loss by using the I3D network \citep{i3d} as feature extractor, and GAN loss with StyleGAN architecture \citep{stylegan}.
% As training objective, we use the following
% \begin{equation}
%     L = L_{VQ} + 0.1L_{Adv} + 0.1L_{\text{Image perceptual}} + 1.0L_{\text{Video perceptual}} + 1.0L_2.
% \end{equation}

% \paragraph{Typical hyper param setting?}
% We need a space for this

% \paragraph{Advantages and Limitations:}
% The ViViT architecture provides a way to encode video into a spatio-temporal latent space that is useful for further processing towards text-to-video generation and has temporal coherence built in.
% However, the ViViT encoder and decoder are limited by the bidirectional attention in time, and it requires to always be given a fixed length video during the encoding process.
% This prevents our text-to-video stage from leveraging a combination of image and video datasets during learning.
% In addition, this limits the inference time by not allowing to optionally provide arbitrary number of inputs frames to condition on for generation, in addition to extending video for arbitrary number of steps by re-applying our method in an autoregressive manner.

% \paragraph{Causal ViViT Architecture:} 
% To overcome the limitations of the ViViT architecture we modify it in the following way. First, the first frame is transformed into a patch independently (no frames are combined). Second, the temporal transformer has a causal mask. The effect of this is that the first frame can be completely independently encoded. This opens up the possibility of text to image training to be embedded naturally into our video model. The second advantage is that we can condition the video generation process on a number of starting frames. 
% \pikinder{Expand this to explain the genius and the novely of this}


%In contrast to the ViViT architecture, the Causal ViViT, or \cvivit in short, can encode videos with number of frames below what is used during training, and also encode single images within the same latent space. \mbz{below or different?} \rubville{less than or equals to} To achieve this, given a video sequence $\bx \in \mathds{R}^{T \times H \times W \times C}$ we tokenize the first frame by linearly projecting non-overlapping $p \times q$ patches into image tokens, and the rest of the frames are tokenized in a similar way as described before into video tokens.
%Once the video tokens are computed, into $\bv \in \mathds{R}^{t \times h \times w \times h}$, the encoding process is the same as in our ViViT encoder, but the attention in the time encoder is causal.
%Therefore, the latent encoding in frame $t$ only observe frames $\leq t$.
%This provides flexible number of input frames and single frame representations so that our text-to-video stage of training can, by design, leverage spatial structure from image datasets and spatio-temporal structure in videos datasets while preserving temporal coherence.

\subsection{Text-to-video generation with bidirectional transformers}
In this stage, the text-to-video task can be formulated as a sequence-to-sequence problem to predict video tokens given the paired text embeddings. Most of recent methods \citep{DALLE, PARTI, nuwa, cogvideo} adopt a  transformer model for these sequence-to-sequence tasks. In their models, they use an auto-regressive transformer which predicts the image or video tokens sequentially given the encoded text features.
% The sampling time scales linearly with the output size and is very slow for high-resolution images and long video sequences.
As a result, the sampling time scales linearly with the sequence length, even when caching is used. This becomes impractical for long video sequence generation. 

\paragraph{Masked bidirectional transformer:}
% \paragraph{Advantages}
% \begin{itemize}
%     \item Inference time
%     \item Quality-time trade off
% \end{itemize}

In this work, we aim to reduce the sampling time by having a small and fixed sampling step disregarding different video sequence lengths. Inspired by previous work for image generation~\citep{MASKGIT}, we use a bidirectional transformer since it can predict different video tokens simultaneously. For training step $i$, we first sample a mask ratio $\gamma_{i}$ from 0 to 1 and randomly replace $\lceil \gamma_{i} \cdot N \rceil$ tokens with the special token \texttt{[MASK]}, where $N$ is the video sequence length. Then
we learn the model parameters by minimizing the cross entropy loss on those masked tokens given the encoded text embeddings and unmasked video tokens. During inference, we first label all of the video tokens as the special token \texttt{[MASK]}. Then, at each inference step, we predict all the masked (unknown) video tokens in parallel conditioned on the text embeddings and unmasked (predicted) video tokens. We keep a ratio $\beta_{i}$ of the predicted tokens at sampling step $i$ and the remaining tokens are re-masked and re-predicted in the next step. 

As discussed in MaskGIT~\citep{MASKGIT}, the masking schedule $\gamma_{i}$ and sampling schedule $\beta_{i}$ have a significant effect on the samples quality therefore we follow the same strategies. Compared to an auto-regressive transformer, the number of sampling steps is an order-of-magnitude smaller (typically we use values in the range of 12 to 48). Generally speaking, more sampling steps improves the quality.

% \subsection{Learning a single model from image and video datasets}
\paragraph{Losses and training strategies:}
% \rubville{Should we use a different notation for the token indices for maskgit training?}
Given a pre-trained \cvivit, videos are encoded into codebook ids $\textbf{a}$ of shape $(t_z + 1)\times w_z \times h_z$ which are flattened into a long vector using the raster ordering from \cite{vit_vqgan}. We then model the text-conditional video token distribution using \textit{Masked Visual Token Modeling} (MVTM)~\cite{MASKGIT}:
\begin{equation}
    L_{\text{mask}}= - \sum\nolimits_{\forall i\in[1,N],m_i=1} \log p(a_i | \textbf{a}_{\bar{M}}, \bp),
\end{equation}
where $\textbf{a}_{\bar{M}}$ represents the masked version of $\textbf{a}$, $m_i$ is a binary variable indicating whether $a_i$ is masked or not, $N$ is the number of video tokens, and $\bp$ is the text condition embedding.
In addition to the MVTM objective, we train using classifier-free guidance by dropping the text condition $10\%$ of the time during training~\cite{cfg,PARTI} .
Finally, we dynamically adjust the MVTM objective during training to allow the use of image and video datasets as a single large dataset.
We achieve this by only applying the masking ratio and objective on the first $w_z \times h_z$ tokens if only a single frame is given or over all video tokens if a full video is given. 
This mixed image and video dataset training strategy allows our models to learn concepts only present in image datasets, and transfer them to concepts present video datasets (e.g., the pencil drawing styled video of the panda in Figure.\ref{fig:compos}).

\paragraph{Inference and auto-regressive generation of long videos:}
At inference time, we sample videos tokens by the same iterative process used in \cite{MASKGIT} with classifier-free guidance scale $\lambda$ to control alignment between the generation and the text condition. Once the first video is generated, we can extrapolate additional frames auto-regressively by encoding the last $K$ generated frames in the last video using \cvivit, initializing MaskGIT with the tokens computed by our \cvivit encoder, and proceed to generate the remaining video tokens conditioned on a text input. During video extrapolation, the text condition can be the same or a different one which enables our model to dynamically create visual transitions between the previous and current text condition visual content, effective generating a visual story an described by the input text.

% \subsection{High fidelity video super resolution}
% In order to increase the fidelity of generated videos, we use video super resolution model BasicVSR~\cite{chan2021basicvsr} and video frame interpolation model CAIN~\cite{choi2020cain}. We train BasicVSR from scratch on a mix of ground truth and generated videos with default hyper-parameters at \url{shorturl.at/ruWY4}. For CAIN, we used the pre-trained checkpoint at \url{shorturl.at/DOZ06}. These models double the spatial and temporal resolution of videos generated by \model.